\section{A public household-finance analysis with unknown causal truth}
\label{s5:public-finance}
Public data expose choices that a known-truth simulation cannot. The archived
401(k) application extract contains 9,915 observations from the SIPP 1991
design used in the DML literature \citep[Section 6.2]{dml2018,public401k}.
The outcome is net financial assets; eligibility and participation are
separate binary variables. The declared controls are age, income, education,
family size, marriage, two-earner status, defined-benefit pension status,
IRA participation and home ownership. All required fields are observed in
the retrieved extract, so no row is removed by the declared missingness rule.
The dollar values are those in the archive, without an inflation conversion.

The first question is the average effect of eligibility on assets under
conditional exchangeability given these controls. This is an identifying
assumption about employer/program selection and saving behavior, not a
consequence of fitting a flexible model. The second question is participation's
complier effect with eligibility as an instrument. It adds exclusion,
instrument independence and monotonicity. The ratio of the adjusted
eligibility effect on assets to its adjusted effect on participation yields
that target under those restrictions. Ordinary participation regression
would answer neither question without a different selection argument.
The graph and observation choices must also establish which controls are
appropriate pre-exposure covariates; the cross-sectional extract alone cannot
verify their timing or rule out unmeasured saving preference.

Two nuisance families are compared: standardized quadratic features with
ridge outcome regression and logistic propensity regression, and histogram
gradient boosting for both roles. Three folds and seeds 7101--7103 are fixed
in advance. The original paper's 100-split aggregation and learner collection
are not replicated numerically; this reproduces the design and score with
declared alternative learners. The matched DoubleML check isolates score
agreement from those learner differences.

\begin{center}\small
\begin{tabular}{llrr}\toprule
Nuisance family&Seed&Eligibility ATE&Participation LATE\\\midrule
Quadratic ridge/logit&7101&8,223&11,905\\
&7102&8,391&12,089\\
&7103&8,932&12,935\\
Histogram boosting&7101&7,968&11,398\\
&7102&8,502&12,261\\
&7103&8,299&11,879\\\bottomrule
\end{tabular}\end{center}
These estimates use archived dollars, with the LATE interpreted per
complier under its additional assumptions. For the first quadratic run,
the eligibility 95\% interval is approximately $[5,737,10,710]$; for the
first boosting run it is $[5,683,10,253]$. First-stage estimates range from
.691 to .699. The three quadratic fits each clip one propensity; the
boosting fits clip none. The largest realized inverse-action weight ranges
from about 32 to 47. Small clipping counts do not prove adequate population
overlap. The raw eligibility mean difference is about 19,559, substantially
larger than the adjusted estimates, but that comparison does not reveal the
unknown causal truth.

The reported covariance uses an iid household working law. The extract does
not supply survey weights, strata or primary sampling units. Consequently
these intervals are a reproducible working analysis, not a design-valid
survey inference claim. Alternative control timing, employer selection,
survey design and weak support require sensitivity or additional data.
The result does not transfer automatically to a bank offer, a different
country or current wealth. The useful product behavior is to preserve the
eligibility/participation target distinction, identical matched-reference
inputs, overlap diagnostics, empirical uncertainty and unverified design
assumptions through the final result.
